Diagramme:
=> Alle Aufgaben gemittelt

•	was mache ich mit dem Reihenfolgeneffekt?

-	Zeit + Genauigkeit pro Zeit
o	am schönsten alles richtig und schnell / effizient
o	falsche Ergbenisse weil verzählt / falsch abgelesen
o	schnell ja, aber zu genauigkeit verhält es sich dann xy
 es geht schneller und nicht signifikant mehr Fehler?? 
 warum machen sie auf dem neuen mehr Fehler?
-> warum kann das so sein?
-> aufgaben falsch gelesen, Überlesen (Fragen mussten so kompliziert sein)
-> beschreiben was ich in den Videos gesehen habe
-> was ist meine Empfehlung (potenzielle neue Arbeiten) die ich gebe, damit diese Forschung in der Zukunft eingesetzt werden kann
 Erwartungskonformität: spiegelt sich in Fehlerraten wieder

-	Zufriedenheit Prümper
o	Fragen aus der endgültigen Bewertung rausgenommen, weil zu konfus
-> nach Usability-Test das Dashboard angepasst, damit die Leute gut damit arbeiten können
o	nur die erste Zeile nehmen für die Ergebnisse (alle, ohne Reihenfolge)
 Schwächen gefunden auch abgedeckt durch den Usability-Test

-	Fehlerbetrachtung – Umfrage
o	in der Zeile verrutscht  fehlende Werte
o	Frage: 2&3
	mit Boxplots (s. \autoref{fig:pruemper_all}) / Histogrammen (s. \autoref{fig:hist_prümper_all}) zeigen, dass 2 Werte ausweichen
	wenn man die rausnimmt ist es signifikant
	sehr hohe Abweichung - Extremwerte
-	Fehlerbetrachtung - Test
o	falsch gelesen (Formulierung kompliziert)
o	Ergebnis geraten (zufällig richtig / falsch)
o	2 Ergebnisse gefragt – nicht richtig gelesen
 Ergebnisse vertauscht
 kein Einfluss auf die Zeiten / Richtigkeit der Ergebnisse
o	es gibt die Unterschiede zu Gunsten des Dashboards
 es gab ein Problem mit der Aufgabe (wurde unterschiedlich interpretiert)
 folgende Probleme wurden mit einbezogen

-	Potenziell am Ende in Methodenkritik den Vergleich der Statistik ziehen – es gibt / gibt keinen Reihenfolgeneffekt – deswegen war smart bei Hälfte der Leute, die Reihenfolge zu tauschen

o	Diskussion
	Fragen nach Norm-Kriterien sortieren
 entsprechend je nach Ergebnis interpretieren was passiert
	Empfehlungen ableiten, was beim Übergang alt-neu wichtig ist
 meine Fragen anhand der Daten beantworten (da reichen die Grafiken um zu sagen, „ja ist besser“ oder „nein ist es nicht“ – in einer Grafik aufzeigen)
-> nach Usability-Test das Dashboard angepasst, damit die Leute gut damit arbeiten können

Notizen Kapitelübersicht
 kann auf eine Seite – aber auch auf 3
-	Ist-Zustand, warum muss was geändert werden (nach(?) Stand der Technik / Stand der Forschung)
-	Darstellung was habe ich Programmiert – was habe ich geändert (Dashboard)
-	Darstellung des Experimentes (Ergebnisse)
